% Transcription — fonds Alexandre Grothendieck, shelfmark 87, batch 2 (pages 21-40).
% Established from the facsimile archives/batches/87/batch-02.pdf (a mirror of
% https://grothendieck.umontpellier.fr/87.pdf; no cover sheet, PDF page k =
% folder page 20+k), per the skill transcribe-grothendieck, with the 700 dpi
% tiles of archives/tiles/87/ and 230-450 dpi re-crops.
%
% Nineteen of the twenty leaves are transcribed. Page 31 is skipped: it carries
% nothing in his hand (only page 32 showing through the sheet), and the gap in
% the page numbers is the record. There is no typescript and nobody else's text
% in the batch. No pagination of his own is visible.
%
% What the batch is about. It continues batch 1's « géométrie des convexes »
% and runs from the cube to the combinatorial polyhedron. Pages 21-22: the
% category of « épures de n-cubes » — a set B of 2n elements with a free
% involution (a relation R with orbits of two) — is equivalent to that of
% k-modules E split into free rank-one summands E_i each carrying a pair
% B_i = -B_i of generators, via E = coprod B_i wedge_{+-1} k; then omega_E ~=
% omega_I wedge bigwedge B_i. Pages 23-26: the combinatorial polyhedron of a
% convex polyhedron (facets ordered by F in closure(F')); the facets of the
% cube C(B,R) are the partial sections of B over I = B/R, order-reversed, and
% an exercise recovers B -> I from the ordered set of partial sections; the
% real cube C_n, Aut_aff(R^n, C_n) = Aut(B_n, sigma_n), so convex n-cubes and
% (oriented) épures of n-cubes are equivalent categories. Pages 27-29: facets
% of a convex set via the relation « x, y interior to one open segment », their
% properties for a polyhedral set (finite, unique maximal facet, closures,
% dimension drop by one, faces of intervals Pi_{F',F''} are again Fac(K)),
% then a first axiomatic definition of a combinatorial polyhedron as a finite
% lattice with the diamond condition. Page 30 counts the low dimensions. Pages
% 32-36: the dimension-2 case as a « carte lisse combinatoire » X_0 = S,
% X_1 = A, X_2 = F with the flags A^{L->} = vec(A) x_A A^{up}, the involutions
% sigma_1, sigma_2, sigma, epsilon, the oriented objects S~, A~, F~ and a
% diagram of the projections; the lattice axiom D is unpacked into d_1, d_2,
% d'_1, d'_2 (p. 34), a lemma on Sup/Inf in an ordered set (p. 35), and
% corollaries: every vertex has at least three edges and every face at least
% three sides, faces are sups of their vertices (p. 36). Pages 37-38: for a
% finite ordered set whose lower intervals realise to spheres (condition R),
% delta(x) = dim |X_{<=x}| is strictly increasing, jumps by one exactly at
% successors, and open intervals realise to spheres of dimension
% delta(y)-delta(x)-2; a lemma for |Y| a topological d-manifold. Pages 39-40:
% the diamond condition (Sigma) defines sigma_{x,z}(y); n-repères x_0 < ... <
% x_n by successors, R_n -> R_{n-1}, and the involutions sigma_0 ... sigma_{n-1}
% with the Coxeter-type commutations; the upper half of page 40 is cancelled.
%
% Notation fixed for the batch: the cursive X of pages 32-40 is \mathfrak{X}
% (noted at first use); the flag set he writes A with a corner-and-arrow
% exponent is A^{\llcorner\!\to}; the set of arcs « tangents » is \vec{A}, the
% « transverses » A^{\uparrow}. The d-list of page 34 follows the page's own
% labels d_1), d_2), d'_1), d'_2).
%
% Notable uncertainties: the margin of page 21 (« épure de n-cube », « cube »
% uncertain) and of page 22 (« vérifier »); page 27's cancelled bracket and
% its circled items (5)-(6), crammed obliquely into the lower-left corner and
% mostly \ill; page 29's (2°) « Pi_{F',F''} … Pi_F » as written; page 33's
% label sigma_2 on the second arrow of the diagram; page 34's cancelled
% paragraph on axiom D, read in part; page 39's cancelled two and a half
% lines and its sideways margin; page 40's cancelled upper half.
%
% Pass: Opus 5.5 (claude-opus-5-5), 2026-09-24 — first pass, unchecked against the pages by a human.

\documentclass[11pt,a4paper]{article}
% Le préambule commun à toutes les transcriptions du fonds Grothendieck.
%
% Il tient en une idée : **l'incertitude se compose, elle ne se résout pas.**
% Une transcription qui lisse un mot douteux a détruit la seule chose pour
% laquelle on la faisait — un lecteur doit pouvoir savoir, à chaque endroit,
% ce qui était sur le feuillet et ce qui a été ajouté. Les six macros qui
% suivent sont donc l'appareil critique, et non de la décoration.
%
% Les mêmes commandes sont reconnues par `scripts/render.mjs`, qui produit la
% vue de lecture du site. Une macro ajoutée ici doit l'être aussi là-bas,
% faute de quoi le rendu échouera bruyamment — ce qui est voulu : un rendu
% silencieusement faux est pire qu'un rendu absent.

\ProvidesPackage{grothendieck}[2026/08/08 Transcriptions du fonds Grothendieck]

\usepackage{amsmath,amssymb,amsthm}
\usepackage{mathrsfs}
\usepackage{stmaryrd}
% Les diagrammes commutatifs sont partout dans ce fonds : c'est la langue
% même de la géométrie algébrique de Grothendieck, et une transcription qui
% les rendrait en prose ne transcrirait rien.
\usepackage{tikz-cd}
% Français par défaut — c'est la langue des feuillets — mais l'anglais est
% chargé aussi : la traduction et le résumé sont des documents anglais, et la
% typographie française (espaces avant les deux-points) y serait une faute.
% Ces documents basculent par \selectlanguage{english} en tête de corps.
\usepackage[english,french]{babel}
\usepackage{fontspec}
\usepackage{geometry}
\geometry{a4paper,margin=25mm,marginparwidth=28mm}
\usepackage{marginnote}
\usepackage{xcolor}
% ulem, not soul. `soul`'s \st and \ul are built on a character-by-character
% scan that XeLaTeX's font machinery breaks: the document fails with an
% undefined control sequence rather than degrading. ulem does the same two jobs
% and is Unicode-engine safe. `normalem` stops it hijacking \emph, which the
% transcriptions use for Grothendieck's own emphasis.
\usepackage[normalem]{ulem}
\usepackage{hyperref}
\hypersetup{colorlinks=true,linkcolor=black,urlcolor=black,pdfborder={0 0 0}}

% --- Métadonnées du lot -----------------------------------------------------
% Reprises telles quelles par le script de rendu, qui les affiche en tête de
% la vue de lecture. Elles fixent ce que le fichier prétend couvrir : sans
% elles, une transcription flotte sans point d'attache dans un fonds de seize
% mille feuillets.

\newcommand{\folder}[1]{\gdef\grothFolder{#1}}
\newcommand{\batch}[1]{\gdef\grothBatch{#1}}
\newcommand{\leaves}[2]{\gdef\grothFirst{#1}\gdef\grothLast{#2}}
\newcommand{\foldertitle}[1]{\gdef\grothFoldertitle{#1}}
\gdef\grothFolder{?}\gdef\grothBatch{?}\gdef\grothFirst{?}\gdef\grothLast{?}\gdef\grothFoldertitle{}

% --- Le résumé --------------------------------------------------------------
%
% Il ouvre la lecture modernisée, sous le titre « Résumé », et il est composé
% en petit. Ce n'est pas un ornement : c'est le seul passage du document qui
% ne soit pas des mathématiques, et un lecteur doit pouvoir le reconnaître
% avant de l'avoir lu — puis le sauter s'il connaît déjà le sujet, ou s'y
% arrêter s'il ne le connaît pas. Le filet qui le suit marque la reprise du
% corps.

\newenvironment{resume}
  {\small\setlength{\parskip}{0.45em}}
  {\par\medskip\hrule\medskip}

% --- Mots-clés ---------------------------------------------------------------
%
% En anglais, en clôture du résumé de la lecture modernisée : le vocabulaire
% moderne sous lequel chercher ce que les feuillets construisent. Le manifeste
% du site (`npm run manifest`) les extrait pour étiqueter la cote entière —
% cette ligne est la source unique des tags, il n'y a pas de fichier à côté.
\newcommand{\keywords}[1]{\par\smallskip\noindent
  {\footnotesize\textsc{Keywords} --- \textit{#1}}\par}

% --- L'appareil critique ----------------------------------------------------

% Le numéro de feuillet, en marge. C'est l'ancre qui permet de régler un
% désaccord sur une formule en regardant *un* feuillet plutôt qu'une chemise
% de six cents. Le site s'en sert aussi pour tourner les pages du fac-similé
% quand on fait défiler la transcription.
\newcommand{\leaf}[1]{%
  \marginnote{\footnotesize\color{gray!70}\textsf{#1}}%
  \label{leaf:#1}\ignorespaces}

% Illisible : on ne devine pas. Un mot inventé qui se lit comme les autres est
% le pire résultat possible.
\newcommand{\ill}{\textcolor{red!60!black}{[\,\ldots\,]}}

% Lecture douteuse : le mot est donné, et signalé comme incertain.
\newcommand{\uncertain}[1]{\uline{#1}}

% Ajout de l'éditeur — une lettre manquante, un mot sous-entendu.
\newcommand{\add}[1]{\textcolor{blue!50!black}{[#1]}}

% Biffé par Grothendieck lui-même. Ce qu'il a rayé fait partie du document :
% une rature dit souvent où la pensée a changé de direction.
\newcommand{\struck}[1]{\sout{#1}}

% Note du transcripteur, sur l'état matériel du feuillet.
\newcommand{\note}[1]{\par\smallskip{\small\color{gray!60!black}\textit{[#1]}}\par\smallskip}

% Note marginale de Grothendieck, distincte du corps du texte.
\newcommand{\marginal}[1]{\par\smallskip\noindent\hspace{1em}%
  {\small\color{gray!60!black}#1}\par\smallskip}

% --- Titre ------------------------------------------------------------------

\AtBeginDocument{%
  \begin{center}
    {\large\bfseries Fonds Alexandre Grothendieck --- cote n\textsuperscript{o}~\grothFolder}\\[2pt]
    {\itshape\grothFoldertitle}\\[4pt]
    {\small feuillets \grothFirst--\grothLast\ (lot \grothBatch)}\\[2pt]
    {\footnotesize\color{gray!60!black}Transcription établie d'après le fac-similé numérisé
     par l'Université de Montpellier.\\
     Les lectures incertaines sont soulignées, les illisibles notées [\,\ldots\,],
     les ajouts entre crochets.}
  \end{center}
  \vspace{1em}}

\endinput


\folder{87}
\batch{2}
\pages{21}{40}
\dating{s.d. — le groupe « Géométrie et topologie combinatoire » (69 à 102) est daté 1976-[vers 1986]}
\watermark{Édition de démonstration}
\foldertitle{Polyèdres convexes : notes manuscrites (s.d.)}

\begin{document}

\page{21}

Soit $n$ un entier, $k$ un anneau \uncertain{comm.}\note{le mot est récrit
au-dessus de la ligne} où $2$ inv. ds $k$.

$C_n$ catégorie formée des

\begin{itemize}
  \item[] $k$-module $E$
  \item[] \struck{\ill{}}\note{un « 1 » souligné est écrit au-dessus du mot
    biffé} $(E_i)_{i\in I}$ de sous-modules libres de $E$
  \item[] $\forall\, i \in I$, partie $B_i$ de $E_i$ ;
\end{itemize}

avec axiomes

\begin{itemize}
  \item[a)] $E$ somme directe des $E_i$ ;
  \item[b)] $\operatorname{card} B_i = 2$ — ($\forall\, i \in I$)
  \item[c)] $\forall\, b \in B_i$, $\{b\}$ base de $E_i$
  \item[d)] $B_i = -B_i$.
\end{itemize}

$\mathcal{D}_n$ catégorie formée des ens. $B$ de card. $2n$, munis d'une
relation d'équiv. $R$ à orbites de cardinal $2$ (donc $B/R = I$ de card. $n$)
\marginal{épure \struck{\ill{}} de $n$-\uncertain{cube} \struck{combinatoire}} — ou encore des ens. de card. $2n$, avec
\struck{\ill{}} $\{\pm 1\}$ opérant librement i.e. avec une involution sans pts
fixes.

\[
  C_n \longrightarrow \mathcal{D}_n \qquad \text{par} \quad
  \bigl(E, (E_i), (B_i)_{i\in I}\bigr) \longmapsto \Bigl\{\textstyle\bigcup B_i, \ldots\Bigr\}
\]

Mais on a aussi
\[
  \mathcal{D}_n \longrightarrow C_n \quad \text{en prenant} \quad
  E = \coprod_{i\in I} B_i \wedge_{\{\pm 1\}} k
\]
où $\{\pm 1\}$ opère sur $k$ par homothéties. Donc les deux catégories sont
équivalentes.

\marginal{$0 \to E \to k^{B} \to k^{I} \to 0$}\note{entouré, à gauche de la formule}
\marginal{ou aussi $E \subset k^{(B)}$ formé des $\sum_{b\in B} \lambda_b e_b$
avec $\lambda_{-b} = -\lambda_b$ \ill{}}

\page{22}

En fait, le $E$ ainsi \uncertain{associé} a \struck{nécessairement} une forme
quadratique canonique, sur chaque $B_i \wedge_{\{\pm 1\}} k$ : la forme
\uncertain{bilinéaire} \add{\uncertain{de produit scalaire}} $(x,y) = xy$.

\uncertain{Je dis} qu'on a un isom. canonique
\[
  \omega_E \simeq \omega_I \wedge \bigwedge_{i\in I} B_i .
\]
En effet, par la formule d'associativité on a
\[
  \omega_E \simeq \omega_I \wedge \bigwedge_{i\in I} \omega_{E_i}
\]
et d'autre part
\[
  \omega_{E_i} \simeq \struck{\ill{}}\, B_i .
\]

\marginal{\ill{} \uncertain{vérifier} :}[On a en effet, comme $B_i \subset$
ens. des bases \add{sur $\mathbb{R}$}\note{un petit « $R$ » est écrit sous la
ligne après « bases »} de $E_i$, $B_i \to \omega_{E_i}$, et on voit que c'est inj,
donc bijectif ……]

\marginal{épure de $n$-cube \struck{Cube combinatoire} orienté}\note{à droite,
en bas de la page, souligné comme un titre}

\page{23}

\emph{Polyèdre combinatoire} associé à un polyèdre convexe : l'ens. des
facettes \add{$\neq$ l'unique facette \uncertain{maxim.} ?}\note{ajout
interlinéaire}, \struck{\ill{}} ordonné par $F \subset \overline{F'}$.

[La dim d'une facette se retrouve sur le polyèdre combinatoire comme une dim.
combinatoire (d'un él. d'un ens. ordonné) : sup des longueurs des chaînes
\uncertain{croissantes}
\[
  F_0 \lneq F_1 \lneq \cdots \lneq F_n = F
\]
— ou longueur commune des chaînes maximales …]

\[
  \Phi_0 \quad \Phi_1 \quad \ldots \quad \Phi_{n-1}\struck{\ill{}}
  \qquad (n = \dim C)
\]
où $\Phi_i$ = ens. des facettes de dim. $i$.

Entre $\Phi_i$ et $\Phi_{i+1}$ ($0 \leq i \leq n$) relation d'incidence $R_i$.

La relation d'ordre sur $\Phi = \coprod_{0\leq i\leq n} \Phi_i$ se reconstitue à
partir des $R_i$.

\emph{Proposition} Les facettes $\neq \emptyset$ du \struck{\ill{}} cube
$C(B,\mathcal{R})$ sont en corr. biunivoque avec l'ens. des sections
partielles de $\mathcal{B}$ sur $I = B/\mathcal{R}$. Si
$\sigma : J \subset I \to \mathcal{B}$ est une telle section, la facette
\emph{fermée} (ouverte) associée est l'ens. des
\[
  \sum_{b\in B} \lambda_b . b \quad \text{avec} \quad
  \lambda_{-b} = -\lambda_b, \quad \lambda_{\sigma(j)} = +1, \quad \text{et}
\]
$|\lambda_b| \leq 1$ (resp. $|\lambda_b| < 1$) pour $b \in I \setminus J$.

\page{24}

Soient $\lambda = \sum_{b\in B} \lambda_b e_b$, $\mu = \sum_{b\in B} \mu_b e_b$
deux éléments de $C$ ($\lambda_{-b} = -\lambda_b$, $|\lambda_b| \leq 1$, idem
pour $\mu$) quand a-t-on $R(\lambda,\mu)$ (relation d'équiv. des facettes ?)

Si $\lambda \neq \mu$, on veut que \struck{\ill{}}
$\exists\, t \in \mathbb{R}$, $t\lambda + (1-t)\mu \in C$, avec au choix $t < 0$,
ou $t > 1$. \struck{\ill{}} On veut \struck{$t\lambda_b$} donc
$|t\lambda_b + (1-t)\mu_b| \leq 1$. Si $\lambda_b = \mu_b$, aucun pb, si
$|\lambda_b|, |\mu_b| < 1$, non plus ; le seul pb si $\lambda_b \neq \mu_b$,
et $\lambda_b$ ou $\mu_b = \pm 1$.

Il faut donc $|\lambda_b| = 1 \Leftrightarrow |\mu_b| = 1$
\[
  \underset{-1}{\bullet} \!\!\!\longrightarrow\!\!\!
  \underset{\lambda_b}{\bullet} \!\!\!\longrightarrow\!\!\!
  \underset{+1 = \mu_b}{\bullet}
  \qquad \Longrightarrow \qquad \lambda_b = \mu_b \quad \text{donc}
\]
\note{le croquis est un segment de $-1$ à $+1$ portant $\lambda_b$ à
l'intérieur et $\mu_b$ à l'extrémité $+1$ ; il est rendu ici par des points
alignés}

Donc \struck{a) L'ens. $J(\lambda)$ des $\lambda$}
\begin{itemize}
  \item[a)] Soit $J(\lambda) = \{\, b \in B \mid \lambda_b = \pm 1 \,\}$. On veut
    $J(\lambda) = J(\mu)$, soit $J$
  \item[b)] Si $j \in J$, on veut $\lambda_b = \mu_b$
\end{itemize}

La proposition s'ensuit.

\emph{Corollaire} La relation d'ordre entre facettes de $C = C(B,R)$ est
inverse de la relation de prolongement entre sections partielles de $B$ sur
$I = B/\mathcal{R}$.

La facette maximale correspond à la section partielle vide ($J = \emptyset$).
Les facettes sous-maximales correspondent aux éléments de $B$. Deux facettes
sous-maximales distinctes \uncertain{ne}\note{ajout interlinéaire} se
rencontrent pas (« sont $/\!/$ ») ssi elles correspondent à des élts $x, y$ de
$B$ qui sont conjugués par $R$.

\page{25}

\emph{Exercice} Soit $B \xrightarrow{p} I$ une application surjective\note{« surjective » est écrit au-dessus de la ligne}
\struck{telle que $\forall\, i \in I$, $\operatorname{card} p^{-1}(i) \geq 2$}.
Soit $\mathcal{P}$ l'ens. des sections partielles définies de $B$ sur (des
parties de) $I$, \struck{\ill{}} ordonné par la relation de prolongement.
\struck{Montrons que les} Ainsi l'ens. des él. maximaux \struck{et \ill{}} de
$\mathcal{P}$ s'identifie à
\[
  C = \prod_{i\in I} B_i \qquad \bigl(B_i = p^{-1}(i)\bigr)
\]
l'ens. des éléments sous-minimaux\note{récrit au-dessus d'un mot biffé}
s'identifie à $B$. Pour tt $b \in B$, soit $C_b \subset C$ l'ens. des él.
maximaux \struck{tels que} qui majorent $b$, i.e. l'ens. des sections $\sigma$
de $B$ sur $I$ telles que $\sigma(p(b)) = b$. Soient $b, b' \in B$,
$b \neq b'$, alors
\[
  C_b \cap C_{b'} = \emptyset \iff p(b) = p(b') .
\]
Donc si $b, b' \in B$, on a
\[
  \bigl(p(b) = p(b')\bigr) \iff \bigl(b = b' \ \text{ou}\ C_b \cap C_{b'} = \emptyset\bigr)
\]
Cela nous permet de reconstruire $B$ et la relation d'équivalence $R$ définie
par $p$ \struck{\ill{}} (donc $I = B/R$ et $p$), à isom. près, en termes de
$\mathcal{P}$.

\emph{Cor.} La cat. des $B \xrightarrow{p} I$ ($p$ surjectif) (pour les isos)
équivalente à une s-catégorie pleine de celle des ens. ordonnés en associant :
à $p$ \ill{} $\mathcal{P}(p)$ ci-dessus.

\page{26}

\subsection*{Cube}

(A)\note{la lettre est entourée} \emph{Pt de vue géométrie des} \struck{\ill{}} \emph{convexes}
(sur $\mathbb{R}$)

\begin{itemize}
  \item[(a)] Cube standard (réel) :
    \[
      C_n \subset \mathbb{R}^n, \qquad
      C_n = \bigl\{\, x = (x_i)_{1\leq i\leq n} \in \mathbb{R}^n \bigm| |x_i| \leq 1 \,\bigr\}
    \]
  \item[(b)] Cube $C$ d'un \struck{vectoriel} espace affine réel $E$ (de dim
    finie $n$) : partie telle qu'il existe isom. \struck{aff} $u : \mathbb{R}^n
    \xrightarrow{\sim} E$ avec $u(C_n) = C$.
  \item[(c)] Catégorie des $n$-cubes \struck{de dim $n$} convexes (dans
    \struck{vectoriels} $\mathbb{R}$-espaces affines de dim. $n$) : c'est une
    groupoïde connexe avec objet marqué $(\mathbb{R}^n, C_n)$, donc équivalent à
    la catégorie des torseurs sous
    \[
      W_n = \operatorname{Aut}_{\text{affin.}}(\mathbb{R}^n, C_n)
    \]
\end{itemize}

\marginal{explicitons en termes de $B \xrightarrow{p} I = B/\mathcal{R}$ :
$E = \operatorname{Hom}_{\{\pm 1\}}(B, \mathbb{R})$,
$C = \{\, \lambda \in \operatorname{Hom}_{\{\pm 1\}}(B, \mathbb{R}) \mid
|\lambda(b)| \leq 1 \ \forall\, b \in B \,\}$}\note{encadré, écrit en long dans
la marge gauche, en regard de la proposition}\emph{Prop} Soit
$B_n = \{\, \pm e_i \mid 1 \leq i \leq n \,\} \subset \mathbb{R}^n$. Alors
\[
  \operatorname{Aut}_{\text{aff}}(\mathbb{R}^n, C_n) = \operatorname{Aut}_{\text{vect}}(\mathbb{R}^n, B_n)
  \xrightarrow{\ \sim\ } \operatorname{Aut}(B_n, \uncertain{\sigma_n})
\]
où $\sigma_n$ est l'involution $b \mapsto -b$ de $B_n$. \struck{relat}

\struck{Dém} Dém. laissée comme exercice. [pas évident, mais pas dur : fait
$\subset$ ; on peut noter que $B_n$ est l'ens. des \emph{pts extrémaux} du
convexe $C_n$ …]

\emph{Corollaire} Catégorie des $n$-cubes convexes équivalente à la catégorie
des épures de\note{« épures de » est écrit au-dessus de la ligne} $n$-cubes
\struck{combinatoires}.

\emph{Corollaire} Cat. des $n$-cubes orientés $\approx$ cat. des épures\note{même
ajout interlinéaire} $n$-cubes \struck{combinatoires} orientés.

\page{27}

\emph{Facettes} d'un ens. convexe $C$ (dans vectoriel réel) : partie
\struck{\ill{}} non vide $F$ de $C$ telle que
\struck{[$x \in F$, $y \in C - F$ implique que $x$ \add{ou $y$} soit une
extrémité de $C \cap D_{xy}$ ($D_{xy}$ = la droite qui passe par $x$ et $y$,
i.e. l'ens. des $tx + (1-t)y$ ($t \in \mathbb{R}$))]}\note{le passage entre
crochets est encadré et barré de traits obliques ; les deux ajouts
interlinéaires qui le suivent, « ou $y \in F$ si » et « ou $y$ », se
rapportent à la même rédaction} \struck{\ill{} convexe} : pour $x \in F$,
$y \in C - \{x\}$, \struck{\ill{}} si $x$ \struck{\ill{}} est une extrémité de
$C \cap D_{xy}$ (i.e. s'il existe $x', y' \in C \cap D_{xy}$ tels que
\[
  \underset{x'}{\bullet}\!-\!\underset{x}{\bullet}\longrightarrow y \longrightarrow y'
\]
\struck{$x, y$ soient} $x, y$ soient \add{intérieurs} entre $x'$ et $y'$) alors
$y \in F$. On peut l'exprimer aussi en disant que tt intervalle ouvert
$\subset C$ qui rencontre $F$ est contenu ds $F$.

\emph{Prop.} Toute facette de $C$ est convexe. $C$ est réunion disjointe de ses
facettes \add{(non vides)}.

Considérons la relation dans $C$
\[
  R(x,y) : \ (x = y) \ \text{\emph{ou}} \ \bigl(x \text{ ni } y \text{ ne sont extrémités de } C \cap D_{xy}\bigr)
\]
i.e. $x, y \in$ \struck{\ill{}} intervalle ouvert $\subset C$ ; c'est une
relation d'équivalence : \struck{symétrie} \uncertain{réflexivité et} symétrie
évident\supplied{s} ; transitivité facile. Il en découle aussi que

\[
  \underset{x'}{\bullet}\,\underset{x}{\bullet} \ \ldots \ \underset{y''}{\bullet}\,\underset{y}{\bullet}\,\underset{y'}{\bullet} \ \ldots \ \underset{z}{\bullet}\,\underset{z'}{\bullet}
\]
\note{le croquis est un triangle de sommets approximatifs $x', y', z'$,
chaque sommet précédé d'un point intérieur $x, y, z$, avec $y''$ près de $y$ ;
il illustre la transitivité et n'est donné ici que par ses étiquettes}

$\forall\, x \in C$, l'ens. des $y \in C$ tels que $R(x,y)$\note{restitution
approximative : la ligne commence par un mot biffé, puis « $x \in C$ »} est
l'unique facette qui contient $x$.

\marginal{ou encore : int. \emph{finie} de demi-espaces fermés …}
\emph{Prop.} \struck{\ill{}} Soit $C$ un \emph{polyèdre} convexe (i.e. un ens.
convexe qui est réunion finie d'ensembles \struck{\ill{}} convexes de partie
finie de l'espace affine $E$) alors : (1)\note{les numéros de cette proposition sont entourés sur la page} l'ens. des facettes est
fini. (2) $\exists$ unique facette maximale $\Phi$, qui est
l'intérieur de $C$ dans \uncertain{son espace} \uncertain{engendré par} $C$.
(3) Pour tt facette $F$, $\overline{F}$ est une réunion de
facettes, et $F$ est l'unique facette $\subset \overline{F}$ de dimension max
$(= \dim \overline{F})$ \struck{\ill{} $\overline{E} = \overline{F}$
$(F, F'$ facettes$)$ … $F = F'$} (4) Si $\dim F = d$,
$\exists$ facette $\subset \overline{F}$ de dim $d-1$ ; si $F, F'$ sont des
facettes, $\overline{F} = \overline{F'} \Rightarrow F = F'$
\marginal{(5) $C$ est l'env. convexe de la réunion des facettes
minimales, i.e. \uncertain{des} pts extrémaux (6) \ill{} bar. des
\ill{} l'ens. des \ill{} \uncertain{caractérise} \ill{}}\note{les deux articles
(5) et (6) sont écrits en oblique dans le coin
inférieur gauche, serrés, et se lisent mal au-delà de leurs premiers mots}
Relation d'ordre entre facettes

\page{28}

$C \subset E$, $C$ partie polyédrale de l'espace affine $E$,
$\dim C = d$

$A_C = \operatorname{Aff}(C)$ espace affine engendré par $C$

$C^0 = \operatorname{Int}_{A_C}(C)$ est la facette maximale (de dim maximale),
$C = \overline{C^0}$.

Les autres facettes \add{fermées} \struck{sont} les intersections \add{$H \cap C$
avec $C$} d'hyperplans $H$ de $A_C$ ne rencontrant pas $C^0$ \struck{\ill{}}.
Les facettes \struck{fermées} \struck{telles que} de dim $d-1$ correspondent
\add{1-1} aux $H$ tels que l'intérieur de $H \cap C$ dans $H$ soit
$\neq \emptyset$, et les intérieurs en question sont les facettes
\add{ouvertes} de dim $d-1$. Soit $C^1$ leur réunion.

Les facettes fermées de dim $\leq d-2$ sont les int. avec $C$ de
sous-espaces affines de dim $d-2$ de $A_C$ ne rencontrant pas
$C^0 \cup C^1 \ldots$

Les facettes fermées de dim $\leq d-i$ sont les int. \add{$F \cap C$} avec
$C$ des sous-espaces affines $F$ de $A_C$ de dim $d-i$ \add{(i.e. codim. $i$)}
\struck{\ill{}} ne rencontrant pas $C^0 \cup C^1 \cup \cdots \cup C^{i-1}$.
Les facettes fermées de dim $d-i$ sont en corr. 1-1 avec les $F$ tels que
$\operatorname{int}_F F \cap C \neq \emptyset$, et les
$\operatorname{int}_F F \cap C$ sont les facettes de dim $d-i$.

Intersection de deux facettes fermées est une facette fermée [c'est une partie
convexe de $C$ et c'est une réunion de facettes, il reste à prouver).

\emph{Prop} Les facettes \add{fermées} de $C$ sont les réunions de facettes qui
sont convexes et fermées.

Soit

\page{29}

$F$ facette, alors \struck{\ill{}}
\begin{itemize}
  \item[a)] $F = \operatorname{int}_{A_F} \overline{F}$ \ ($A_F = A_{\overline{F}}$),
    \ $\overline{F} = A_F \cap C$ = \struck{\ill{} facettes},
  \item[b)] $F' \subset C$ alors
    $\underbrace{F' \text{ facette de } F}_{\Updownarrow}$ \
    $F'$ facette de $C$, $F' \subset F$ \
    [$\Rightarrow \overline{F}$ réunion de facettes]
  \item[c)] Soit \struck{\ill{}} $K \subset C$. Pour que $K$ soit une facette
    fermée, il faut et s. que $K$ soit convexe, fermé et une réunion de
    facettes. ($\Rightarrow$ une intersection de deux facettes est une facette)
  \item[d)] \struck{\ill{}} Soit $\operatorname{Fac}(C) = \Pi$, \add{l'ens.
    ordonné des facettes}. Alors
\end{itemize}

\begin{itemize}
  \item[1°)] $\Pi^{\circ}$ (l'ens. ordonné opposé) est
    \uncertain{isomorphe} à un ens. $\operatorname{Fac}(C')$, où $C'$ est
    une \add{partie polyédrale} \struck{convexe} \struck{\ill{}} dans $E$ de
    $=$ dim.
  \item[2°)] Soient $F', F'' \in \Pi$, \add{avec $F' < F''$},
    $\Pi_{F',F''}$ l'ens. ordonné des $F \in \Pi$ tels que
    $F' \leq F \leq F''$. Alors $\Pi_F$ est \add{$\simeq$ à un ens.
    ordonné} de la forme \struck{\ill{}} $\operatorname{Fac}(K)$, où $K$ un
    \add{polyèdre} convexe de dim $\delta = \dim F'' - \dim F' - 1$
\end{itemize}

\emph{Cor} Si \struck{\ill{}} $\dim F'' - \dim F' = 2$, alors
\struck{card} l'ens. des $F$ tels que $F' < F < F''$ est de card. $2$.

\marginal{\ill{}}\note{à gauche, un petit croquis : quatre segments issus d'un
point, dans un cercle, deux d'entre eux repassés à l'encre}
\emph{Polyèdre combinatoire} : ens. ordonné fini $I$ satisfaisant les
conditions
\begin{itemize}
  \item[a)] $I$ a des inf et des sup finis [plus petit et plus grand élément,
    et $i \wedge j$ et $i \vee j$]
  \item[b)] $I$ satisfait condition des chaînes
  \item[c)] Si $i < j < k$ et si $j$ succ. \struck{immédiat} de $i$, $k$ succ.
    de $j$
\end{itemize}

\page{30}

alors $\exists\, j' \in I$ tel que $j' \neq j$, $j'$ succ. de $i$, $k$ succ. de
$j'$

\begin{itemize}
  \item[$\Rightarrow$ d)] \struck{\ill{}} $\forall\, i \in I$,
    \add{$i$ non minimal} \struck{\ill{}}
    $i = \sup\,(\, j \in I \mid j < i \,)$ ;
    \add{$i$ non max-imal} $\Rightarrow$ $i = \operatorname{Inf}\,(\, j \in I \mid j > i \,)$
\end{itemize}
\note{la condition d) est reliée par une flèche à la condition c), dont elle
découle ; les deux ajouts « $i$ non minimal », « $i$ non maximal » remplacent
une incise biffée où se lit « plus petit »}

\emph{dim $-1$} — $I$ a un seul élément ($0 = 1$).

\emph{dim $0$} — $I$ a deux éléments distincts ($0$ et $1$, $0 \neq 1$)

\emph{dim $1$} — $I$ a \struck{trois} quatre éléments $0$, $1$ et deux
incomp. $x, y$ \struck{($x \neq y$ \ill{})}

\begin{tikzcd}[column sep=small, row sep=small]
  & 1 & \\
  x \arrow[ur, no head] & & y \arrow[ul, no head] \\
  & 0 \arrow[ul, no head] \arrow[ur, no head] &
\end{tikzcd}

\emph{dim $2$}
\begin{itemize}
  \item[] dim $-1$ — $0$
  \item[] dim $0$ — $x_0 \ x_1 \ldots x_{n-1}$
  \item[] dim $1$
  \item[] dim $2$ — $1$ — contours dont les comp. connexes ont au moins
    $3$ sommets
\end{itemize}
\note{à droite, deux polygones dessinés, un hexagone dont un sommet est
marqué et un pentagone}

\emph{dim $3$} — On voudrait que si $x$ est un sommet les $y > x$,
\struck{\ill{}} \add{$y \neq 1$}, forment un contour \emph{convexe} et si $x$
est une face les $y < x$, $y \neq 0$, forment un \emph{contour}
($\uncertain{\pi_0 = 0}$)
\begin{itemize}
  \item[] $-1$ — $0$
  \item[] $0$ — sommets
  \item[] $1$ — arêtes
  \item[] $2$ — faces
  \item[] $3$ — $1$
\end{itemize}

\emph{dim $4$}

\page{32}

\begin{tikzcd}[column sep=small, row sep=small]
  & \vec{A} \arrow[dl] \arrow[dr] & & A^{\uparrow} \arrow[dl] \arrow[dr] & \\
  S = \mathfrak{X}_0 & & A = \mathfrak{X}_1 & & F = \mathfrak{X}_2
\end{tikzcd}

\note{sur la page, les égalités $S = \mathfrak{X}_0$, $A = \mathfrak{X}_1$,
$F = \mathfrak{X}_2$ sont écrites verticalement sous $S$, $A$, $F$ ; la lettre
rendue $\mathfrak{X}$ est un $X$ cursif, et $A^{\llcorner\!\to}$ ci-dessous
rend un exposant qui réunit la flèche verticale de $A^{\uparrow}$ et la
flèche horizontale de $\vec{A}$}

\begin{itemize}
  \item[a)] $\vec{A} \hookrightarrow S \times A$, \struck{$\forall$}
    $\vec{A} \to A$ de rang $2$
  \item[a')] $A^{\uparrow} \hookrightarrow A \times F$, \ $A^{\uparrow} \to A$
    de rang $2$
  \item[b)] $A^{\llcorner\!\to} = \vec{A} \times_A A^{\uparrow} \to S \times F$
    est de rang $2$ sur son image
  \item[c)] $\forall\, s \in S$, l'ens $\mathfrak{X}_{\geq s}$ des $a \in A$ et
    $f \in F$ incidents à $s$ (qui est un \emph{graphe} de sommets $A$,
    d'arêtes $F$, grâce à b), et un \emph{contour} grâce à a')) est un
    \emph{polygone} (i.e. est un contour \add{$\neq \emptyset$}
    \emph{connexe})
  \item[c')] $\forall\, f \in F$, l'ens $\mathfrak{X}_{\leq f}$ des $a \in A$
    et $s \in S$ incidents à $f$ (qui est un graphe … grâce à a), et un
    \emph{contour} grâce à b)) est un \emph{polygone} (i.e. est un contour
    \add{$\neq \emptyset$} \emph{connexe})
\end{itemize}

Version \uncertain{commune} (en introduisant $\mathfrak{X}_{-1} = \{\emptyset\}$,
$\mathfrak{X}_3 = \{1\}$ (plus grand et plus petit élément))

Si $-1 \leq i < j \leq 3$, \add{avec} \uncertain{$\delta =$} $j - i \geq 2$,
et \struck{\ill{}} $(i,j) \neq (-1,3)$, si $x_i \in \mathfrak{X}_i$,
$y_j \in \mathfrak{X}_j$, et $\mathfrak{X}_{\,]x_i,y_j[}$ est l'ens ordonné
des $x \in \mathfrak{X}$ tels que $x_i < x < y_j$, alors la réalisation
topologique de $\mathfrak{X}_{\,]x_i,y_j[}$ est homéomorphe à
\struck{\ill{}} $S^{\delta-1}$ (sphère de dim $\delta - 1$). [Se généralise aux
polyèdres lisses de dim quelconque, cf plus bas …]
\note{« $\delta =$ » est entouré, et « avec » écrit au-dessus ; l'indice de
$\mathfrak{X}$, deux fois, est lu d'après la définition qui le suit}

Les conditions a), a'), b) permettent de définir sur $A^{\llcorner\!\to}$ les
opérations $\sigma_1, \sigma_2$ et $\varepsilon$ [avec
$\sigma_1\sigma_2 = \sigma_2\sigma_1$ \struck{\ill{}}, soit $\sigma$, et
$\sigma_1, \sigma_2, \varepsilon$ étant des involutions,
$\sigma_1, \sigma_2, \sigma, \varepsilon$ étant ss pts fixes —] les conditions
c), c') assurent que $S$, $F$ « sont » bien les ens de sommets et des faces de
la carte lisse combinatoire ainsi définie …

\page{33}

On trouve ensuite
\begin{itemize}
  \item[] $\widetilde{S} = A^{\llcorner\!\to}/\rho_s \simeq$ ens des couples
    $(s,\omega)$, $s \in S$, $\omega$ une orientation du \struck{\ill{}}
    polygone $\mathfrak{X}_{\geq s}$ des $a \in A$ et $f \in F$ incidents à
    $s$ [i.e. un \add{des 2} ordres circulaires sur l'ens des $a \in A$
    incidents à $s$, qui définissent sa structure polygonale]
  \item[] $\widetilde{F} = A^{\llcorner\!\to}/\rho_f$ : ens des couples
    $(f,\varpi)$, $f \in F$, $\varpi$ une orientation du polygone
    $\mathfrak{X}_{\leq f}$ des \struck{\ill{}} $a \in A$ et $s \in S$
    incidents à $f$, [i.e. un des 2 ordres circ. sur l'ens. des $a \in A$
    incidents à $f$, qui définissent sa structure polygonale]
  \item[] $\widetilde{A} = A^{\llcorner\!\to}/\sigma$ = ens des couples
    $(a,u)$, $a \in A$, $u \in \vec{A}(a) \wedge A^{\uparrow}(a)$, où
    $\vec{A}(a)$ ($A^{\uparrow}(a)$) est l'ens des arcs tang. (ou
    transverses) \uncertain{portés} par l'arête $a$.
\end{itemize}

\[
  A^{\wedge} = \operatorname{Im}(A^{\llcorner\!\to} \to S \times F)
  = \text{rel. d'incidence entre $S$ et $F$}
\]

\begin{tikzcd}[column sep=small, row sep=small]
  & & A^{\llcorner\!\to} \arrow[dll, "\rho_S"'] \arrow[dl, "\uncertain{\sigma_2}"'] \arrow[d, "\sigma"] \arrow[dr, "\sigma_1"] \arrow[drr, "\rho_F"] & & \\
  \widetilde{S} \arrow[d, "\sigma_1"'] & \vec{A} \arrow[dl] \arrow[dr, "\sigma_1"] & \widetilde{A} & A^{\uparrow} \arrow[dl, "\sigma_2"'] \arrow[dr] & \widetilde{F} \arrow[d, "\sigma_2"] \\
  S & & A & & F
\end{tikzcd}

\note{sous $\widetilde{A}$ est écrit « $\sigma_1 \updownarrow \sigma_2$ »,
sans flèche qui le rattache au reste du diagramme}

\begin{tikzcd}[column sep=small, row sep=small]
  & A^{\llcorner\!\to} \arrow[dl] \arrow[d] \arrow[dr] & \\
  \widetilde{S} \arrow[d] & A^{\wedge} \arrow[dl] \arrow[dr] & \widetilde{F} \arrow[d, no head] \\
  S & & F
\end{tikzcd}

\note{le trait de $\widetilde{F}$ à $F$ n'a pas de pointe sur la page}

\begin{align*}
  A^{\llcorner\!\to} &\simeq \widetilde{S} \times_S \vec{A} \simeq \struck{\ill{}}\, A^{\uparrow} \times_F \widetilde{F} \\
  &\simeq \vec{A} \times_A A^{\uparrow} \\
  &\simeq \vec{A} \times_A \widetilde{A} \simeq \widetilde{A} \times_A A^{\uparrow} \\
  &\simeq \widetilde{S} \times_S A^{\wedge} \simeq A^{\wedge} \times_F \widetilde{F}
\end{align*}

$A^{\llcorner\!\to}$ s'identifie aux rel. d'incidence entre $(S, A, F)$
simultanées
\begin{itemize}
  \item[] entre $\widetilde{S}$, $\widetilde{A}$,
  \item[] entre $\widetilde{A}$, $\widetilde{F}$,
  \item[] entre $\widetilde{S}$, $\widetilde{F}$,
\end{itemize}

\page{34}

Explicitons l'axiome

D) Toute \add{partie} \struck{\ill{}} minorée \add{$\neq \emptyset$}
\struck{\ill{}} de
$\mathfrak{X}_{**} = \mathfrak{X}_0 \cup \mathfrak{X}_1 \cup \mathfrak{X}_2$
admet un inf, ($\Leftrightarrow$ toute partie majorée \add{$\neq \emptyset$}
de $\mathfrak{X}_{**}$ admet un sup).

\struck{Comme il résulte de ce que les él. minimaux de $\mathfrak{X}_{**}$
sont les él. de $S$, les parties minorées de $\mathfrak{X}_{**}$ sont
celles contenues dans un \add{(ou plus petit él.)} $\mathfrak{X}_{\geq s}$.
Or comme un ordonné \ill{} est un polygone, et dans un polygone
\uncertain{satisfait} la condition \add{d'existence} de borne inférieure ses
\ill{} \ill{} \ill{} est $\geq 3$ (dans le cas excluant alors $n = 2$
ici …) Donc la condition D se réduit :}\note{ce paragraphe est encadré et
barré de traits obliques}

\struck{d) $\forall\, s \in S$, il y a au moins trois arêtes \add{$a \in A$}
incidentes à $s$.}

Il suffit d'exprimer que deux él. distincts \add{minimaux distincts
\uncertain{majorés}} de $\mathfrak{X}_{**}$ \ill{} ont un inf. Si l'un des
deux est $\in \mathfrak{X}_0$ (i.e. est minimal) c'est trivialement
satisfait. Il reste : 3 \uncertain{cas}
\begin{itemize}
  \item[1°)] \emph{Deux arêtes} : on trouve que 2 arêtes qui ont \ill{}
    sommets sont égales, i.e. $(S, A, \vec{A})$ est un graphe strict
  \item[2°)] \emph{Deux faces} : on trouve \struck{qu'elles \ill{}}
    \struck{que les deux faces distinctes de $(F, A, A^{\uparrow})$} \add{qui
    ont en commun deux sommets}, \struck{elles-mêmes} ceux-ci sont liés
    \struck{\ill{}} par l'arête qui les joint, \add{\ill{} \uncertain{fait}}
    \struck{et l'un d'eux \ill{} sommet}
  \item[3°)] Une arête \add{$a$} et une face \add{$f$} \struck{non} incident\supplied{e}s
    — mais ayant un sommet commun : \struck{il} faut exprimer que le deuxième
    sommet de $a$ n'est pas dans $f$, ou encore \struck{\ill{}} : si $f$ est
    incidente aux deux sommets d'une arête, elle est incidente à l'arête.
\end{itemize}

Regroupons, on trouve

\marginal{toute arête est le sup de l'ens. de ses 2 extrémités}\note{en
regard de d$_1$) et d$_2$), qu'une accolade réunit}
\begin{itemize}
  \item[d$_1$)] Une arête est déterminée par \struck{ses deux} l'ens de ses
    extrémités
  \item[d$_2$)] Une face contenant les 2 extr. d'une arête contient celle-ci
\end{itemize}

\marginal{toute arête est le Inf de l'ens des deux faces incidentes}\note{en
regard de d$'_1$) et d$'_2$), qu'une accolade réunit}
\begin{itemize}
  \item[d$'_1$)] Une arête est déterminée par l'ens des deux faces incidentes
  \item[d$'_2$)] Les deux faces incidentes à une arête n'ont d'autre sommet
    incident commun que ses extrémités
\end{itemize}

\page{35}

[$\mathfrak{X}$ un ordonné, $A \subset \mathfrak{X}$, $B$ l'ens des
$b \in \mathfrak{X}$ qui majorent $A$ \struck{\ill{}} Par définition,
$\operatorname{Sup}_{\mathfrak{X}} A$ existe ssi $B$ a un plus petit élément
$\beta$, et on pose \struck{\ill{}} $\operatorname{Sup}_{\mathfrak{X}} A = \beta$.
Si p. ex. $A$ est vide, $B = \mathfrak{X}$, et
$\operatorname{Sup}_{\mathfrak{X}} \emptyset$ existe ssi $\mathfrak{X}$ a un
plus petit élément (qui est alors $\operatorname{Sup}_{\mathfrak{X}} \emptyset$).
On définit dualement $\operatorname{Inf}_{\mathfrak{X}} A \ldots$

\emph{Prop 1} Si $A$ a un plus grand élément \add{$\alpha$}, il a un sup, et c'est
$\alpha$ (Trivial)\note{la Prop 1 est écrite serrée entre les lignes, au-dessus
de la Prop 2}

\emph{Prop 2} Pour que $\operatorname{Sup}_{\mathfrak{X}} A$ existe, i.e. $B$
ait un plus petit él., il f. et s. que $\operatorname{Inf}_{\mathfrak{X}} B$
existe, et alors $\operatorname{Inf}_{\mathfrak{X}} B = \operatorname{Sup}_{\mathfrak{X}} A$.


C'est évidemment \uncertain{néc.} d'après prop. 1. Suff ? Soit
$\widetilde{A}$ l'ens des $a \in \mathfrak{X}$ qui minorent $B$. Donc
$A \subset \widetilde{A} < B$. Par hypothèse (B admet un Inf) $\widetilde{A}$ a
un plus grand él. $\alpha$, et que majorant de $\widetilde{A}$ et a fortiori
\uncertain{majorant} de $A$, donc $\alpha \in B$, et étant dans
$\widetilde{A}$ il minore $B$, donc il est un plus petit él. de $B$, cqfd.

\emph{Cor} Si tte partie de $\mathfrak{X}$ admet un Inf, alors tte partie de
$\mathfrak{X}$ admet un Sup, et réciproquement. \add{\uncertain{conj.
\ill{}}} Si tte partie \struck{\ill{}} minorée \add{de $\mathfrak{X}$} admet un
Inf, alors tte partie non vide (majorée) de $\mathfrak{X}$ admet un Sup, et
réciproquement. Donc pour que tte partie min. non vide de $\mathfrak{X}$
admette un Inf, il f. et s. que tte partie majorée non vide de
$\mathfrak{X}$ admette un Sup.]

\page{36}

Montrons que l'axiome $D = $ d) + d') implique

\emph{Corollaire 1} (moyennant l'axiome D) pour tt $s \in S$,
\struck{\ill{}} $p(s) = \operatorname{Card}\bigl(A\{s\} = \{\, a \in A \mid a \geq s \,\}\bigr) \geq 3$,
et pour tout $f \in F$,
$q(f) = \operatorname{Card}\bigl(A\{f\} = \{\, a \in A \mid a \leq f \,\}\bigr) \geq 3$.

En effet, \struck{on aurait} $p(s) \neq 0$, en vertu de c). Si on avait
$p(s) = 1$, on aurait par c) une \add{unique} seule face incidente à l'arête
$a$ incidente à $s$, contrairement à b). Si on avait $p(s) = 2$, en vertu de
c) les deux arêtes issues de $s$ seraient incidentes simultanément aux deux
faces incidentes à $s$, contrairement à d$'_1$). Donc $p(s) \geq 3$. La relation
$q(f) \geq 3$ s'établit de même.
\marginal{\emph{NB} $p(s) \geq 2$ est indépendant de D, de même
$q(f) \geq 2$.}

\emph{Corollaire 2} (moyennant l'axiome D) : toute face $f \in F$ est le sup
de l'ens de ses sommets, et le sup de l'ens de ses arêtes. Tout sommet est le
inf de l'ens des faces qui y sont incidentes, et le inf de l'ens des arêtes
incidentes.

\page{37}

Soit $\mathfrak{X}$ un ens. ordonné fini (on ne suppose pas qu'il ait un plus
petit ou plus grand élément). On note \add{$|\mathfrak{X}| =$}
réal.$(\mathfrak{X})$ sa réalisation topologique (avec sa décomposition \ill{}
canonique, permettant d'identifier $\mathfrak{X}$ à un ens. ordonné de parties
fermées de $\mathfrak{X}$ \add{recouvrant $\mathfrak{X}$}). Pour tt
$x \in \mathfrak{X}$, soit
$\mathfrak{X}_{<x} = \struck{\ill{}}\ \{\, y \mid y \in \mathfrak{X},\ y < x \,\}$.
Faisons l'hypothèse :

(R) $\forall\, x \in \mathfrak{X}$, $|\mathfrak{X}_{<x}|$ est une sphère
topologique.

Soit $\delta(x) = \dim |\mathfrak{X}_{<x}| + 1 = \dim |\mathfrak{X}_{\leq x}|$.
On trouve

\emph{Proposition} Si $x < y$, on a $\delta(x) < \delta(y)$ (i.e.
$\delta : \mathfrak{X} \to \mathbb{N}$ est strictement croissante). Pour qu'on
ait \struck{\ill{}} $\delta(y) = \delta(x) + 1$ il f. et s. que $y$ soit un
successeur de $x$.

Première assertion claire, compte tenu que $\dim \underbrace{C(X)}_{\text{cône sur } X} = \dim X + 1$,
et $\dim$ est une fonction croissante … L'assertion « il faut » en
résulte, prouvons « il suffit ». Supposons que $y$ soit un successeur de $x$.
Considérons $S = |\mathfrak{X}_{<y}|$ sphère de dim
$d = \delta(y) - 1 \geq \delta(x)$, contenant $|\mathfrak{X}_{\leq x}|$, qui
est une facette maximale, donc de dim $= d$, donc $\delta(x) = \delta(y) - 1$,
cqfd. [Pour qu'on ait $\delta(y) = \delta(x) + 1$, il suffit que $y$ soit
succ. de $x$ et $|\mathfrak{X}_{<y}|$ une \emph{variété} topologique]
\marginal{vérifier en général pour variété}

\page{38}

\emph{Cor} Moyennant (R), $\mathfrak{X}$ satisfait la condition des chaînes, et
$\delta$ est l'unique fonction $\mathfrak{X} \to \mathbb{N}$ telle que
\begin{itemize}
  \item[a)] $\delta(y) = \delta(x) + 1$ si $y$ est un successeur de $x$
  \item[b)] $\delta(x) = 0$ si $x$ élément minimal de $\mathfrak{X}$.
\end{itemize}

\emph{Prop} \add{(Moyennant R)} Soit $x < y$, alors
$|\mathfrak{X}_{x<\cdot<y}|$ est une sphère topologique de dim.
$\delta(y) - \delta(x) - 2$ (on convient de considérer l'espace vide comme
sphère top. de dim. $-1$)\note{dans la marge gauche, en regard de cette
proposition et du lemme qui suit, deux signes serrés que l'on ne lit pas}

Ceci résulte du

\emph{Lemme} Soit $\mathfrak{Y}$ un ens ordonné fini tel que $|\mathfrak{Y}|$
est une $d$-variété topologique. Alors
\begin{itemize}
  \item[a)] $\forall\, x \in \mathfrak{Y}$, $|\mathfrak{Y}_{>x}|$ est une sphère
    top. de dim $d - \delta(x) - 2$
  \item[b)] $\forall\, x \in \mathfrak{Y}$, $|\mathfrak{Y}_{<x}|$ est une sphère
    top. de dim $\delta(x) - 1$ (i.e. $\mathfrak{Y}$ satisfait la condition
    (R))
  \item[c)] $\forall\, z, x \in \mathfrak{Y}$, $z < x$, $|\mathfrak{Y}_{z<\cdot<x}|$
    \struck{satisfait} est une sphère topologique de dim
    $\delta(x) - \delta(z) - 2$.
\end{itemize}

\page{39}

$\mathfrak{X}$ un ordonné fini. On suppose qu'il satisfait la condition
\add{$\overline{\mathfrak{X}} = \mathfrak{X} \cup$ plus petit él.}

($\Sigma$) Soient $x \in \overline{\mathfrak{X}}$, $y$ un successeur de $x$,
$z$ un successeur de $y$. Alors il existe un unique $y' \in \mathfrak{X}$, tel
que $y' \neq y$, $x < y' < z$ ;
\struck{\ill{} $y'$ est un successeur de $x$ et \ill{} minimal de
$\mathfrak{X}$, $z$ un successeur, $\exists\, y' \in \mathfrak{X}$ tel que
$y' < z$ \ill{} et $y'$ est minimal et $z$ un successeur de $y'$}\note{deux
lignes et demie barrées, avec un ajout interlinéaire lui-même barré, où se lit
« si $y$ est un él. »}

Le $y'$ en question (noté $\sigma_{x,z}(y)$) est un successeur de $x$, et $z$
en est un successeur. En effet, si on avait \add{p. ex.} $x < y'' < y'$, on
aurait $x < y'' < z$, donc (comme $y'' \neq y'$) $y'' = y$ donc
$y < y' < z$, contrairement à l'hyp. que $z$ est un succ. de $y$.

\begin{tikzcd}[column sep=small, row sep=small]
  & y' \arrow[dr] & \\
  x \arrow[ur] \arrow[dr] & & z \\
  & y \arrow[ur] &
\end{tikzcd}

\note{croquis en marge gauche ; au-dessus, un premier croquis du même
losange, raturé}

\marginal{\emph{On suppose} pour $x \in \mathfrak{X}$ les chaînes
\struck{\ill{}} \add{\uncertain{issues de $x$}} \struck{\ill{}} maximales
\struck{\ill{}} toutes de longueur $n$.}\note{écrit en long dans la marge
gauche, en regard de ce qui suit}
Un \emph{repère} de $\mathfrak{X}$ \add{($n \geq 0$)} : suite
$x_0, x_1, \ldots, x_n$ d'éléments de $\mathfrak{X}$, tels que
\begin{itemize}
  \item[a)] $x_0$ pt minimal de $\mathfrak{X}$
  \item[b)] $\forall\, 0 \leq i \leq n-1$, $x_{i+1}$ est un successeur de $x_i$.
\end{itemize}
Soit $R_n(\mathfrak{X})$ l'ens des $n$-repères.

\emph{Structures}
\begin{itemize}
  \item[a)] $R_n \xrightarrow{p_n} R_{n-1}$ — ($n \geq 1$)
  \item[b)] Opérations $\sigma_0, \sigma_1, \ldots, \sigma_{n-1}$ — (si $n \geq 1$)
\end{itemize}

\emph{Propositions}
\begin{itemize}
  \item[a)] Les $\sigma_i$ sont des involutions ($n \geq 1$)
  \item[b)] Soit $n \geq 3$, et $0 \leq i < j \leq n-1$, avec $j \neq i+1$,
    alors $\sigma_i$ et $\sigma_j$ commutent \struck{\ill{}} (dans $R_n$).
  \item[c)] Soit $0 \leq i_1 < i_2 < \cdots < i_\nu \leq n-1$, \add{($\nu \geq 1$)}
    avec $i_{\alpha+1} \neq i_\alpha + 1$ $\forall\, \alpha = 1, \ldots, \nu-1$.
    Alors $\sigma_{i_1}\sigma_{i_2}\cdots\sigma_{i_\nu}$ (qui est une
    involution de $R_n$) est sans pt fixe.
  \item[d)] $R_n \xrightarrow{p_n \cdots p_{p+1}} R_p$ commute à
    $\sigma_{n-1}, \ldots, \sigma_{p+1}$ et $\sigma_{p-1}, \ldots, \sigma_0$
\end{itemize}

\page{40}

\note{la moitié supérieure de la page, jusqu'à « 2 sommets … », est
barrée de trois longs traits obliques ; elle se lit néanmoins et est donnée
ici. En marge, plusieurs petits croquis de sommets et d'arêtes, dont un
losange, un éventail hachuré et un faisceau de flèches, qui ne sont pas
reproduits}

\struck{$\mathfrak{X}_0 = S$, $\mathfrak{X}_1 = A$, $\mathfrak{X}_2 = F$}\note{les égalités sont écrites verticalement, chaque
$\mathfrak{X}_i$ au-dessus de sa lettre}

\marginal{$x_0 \ldots x_p\ x_{p+1} \ldots x_{n-1}\ x_n$}\marginal{$y_1 - y_2 \ldots y_n$}\marginal{$x_0 - x_1 - x_2 \ldots - x_n$}\note{en
haut à droite, trois suites superposées dans cet ordre ; une accolade est
tracée sous $x_{p+1} \ldots$, et la suite des $y_i$ se raccorde par un trait
à $x_0$}

\struck{a) $\forall\, f \in F$, les $s \in S$ et $a \in A$ incidents : $F$ forme
un polygone, donc a) si $a < f$ ($a \in A$) $a$ a exactement 2 sommets b) si
$s < f$ ($s \in S$) il y a exactement 2 arêtes entre $s$ et $f$) c) Condition
de connexité}

\struck{On trouve aussi que les arêtes incidentes à une face (\ill{} ou autour
de $|\mathfrak{X}|$ des pts au voisinage desquels $|\mathfrak{X}|$ est de dim
1 !) et \ill{} on conclut que les arêtes a exactement 2
sommets …}

… en faisant que $\sigma_{n-1}, \ldots, \sigma_{p+1}$ triviales sur $R_p$
$(n > p \geq 0)$\note{cette ligne, non barrée, paraît compléter l'article d)
de la page 39}

\end{document}
